\subsection{SYSTEMS OF VECTORS WITH NEGATIVE SCALAR PRODUCTS}

Lemma 3. Let \(q\) be a positive quadratic form on a real vector space \(V\) and let \(B\) be the associated symmetric bilinear form. Let \(a_1, \ldots, a_n\) be elements of \(V\) such that \(B(a_i, a_j) \leqslant 0\) for \(i \neq j\) .

\begin{enumerate}
\def\labelenumi{(\roman{enumi})}
\tightlist
\item
  If \(c_{1}, \ldots, c_{n}\) are real numbers such that \(q\left(\sum_{i} c_{i} a_{i}\right) = 0\) , then
\end{enumerate}

\[
q \left(\sum_ {i} | c _ {i} | a _ {i}\right) = 0.
\]

\begin{enumerate}
\def\labelenumi{(\roman{enumi})}
\setcounter{enumi}{1}
\tightlist
\item
  If \(q\) is non-degenerate and if there exists a linear form \(f\) on \(V\) such that \(f(a_i) > 0\) for all \(i\) , the vectors \(a_1, \ldots, a_n\) are linearly independent.
\end{enumerate}

The relation \(\mathrm{B}(a_i, a_j) \leqslant 0\) for \(i \neq j\) immediately implies that

\[
q \left(\sum_ {i} | c _ {i} | a _ {i}\right) \leqslant q \left(\sum_ {i} c _ {i} a _ {i}\right),
\]

hence (i). If \(q\) is non-degenerate, the relation \(\sum_{i} c_{i} a_{i} = 0\) thus implies that

\[
\sum_ {i} | c _ {i} | a _ {i} = 0;
\]

it follows that, for any linear form \(f\) on \(\mathbf{V}\) , we have \(\sum_{i}|c_i|f(a_i) = 0\) , and hence \(c_i = 0\) for all \(i\) if we also have \(f(a_i) > 0\) for all \(i\) . This proves (ii).

Lemma 4. Let \(\mathbf{Q} = (q_{ij})\) be a real, symmetric, square matrix of order \(n\) such that:

\begin{enumerate}
\def\labelenumi{\alph{enumi})}
\item
  \(q_{ij} \leqslant 0\) for \(i \neq j\) ;
\item
  there does not exist a partition of \(\{1,2,\ldots,n\}\) into two non-empty subsets I and J such that \((i,j)\in\mathrm{I}\times\mathrm{J}\) implies \(q_{ij}=0\) ;
\item
  the quadratic form \(q(x_{1},\ldots ,x_{n}) = \sum_{i,j}q_{ij}x_{i}x_{j}\) on \(\mathbf{R}^n\) is positive.
\end{enumerate}

Then:

\begin{enumerate}
\def\labelenumi{(\roman{enumi})}
\item
  The kernel \(\mathbf{N}\) of \(q\) is of dimension 0 or 1. If \(\dim \mathbf{N} = 1\) , \(\mathbf{N}\) is generated by a vector all of whose coordinates are \(> 0\) .
\item
  The smallest eigenvalue of \(\mathbf{Q}\) is of multiplicity 1 and a corresponding eigenvector has all its coordinates \(>0\) or all its coordinates \(< 0\) .
\end{enumerate}

Since \(q\) is a positive quadratic form, the kernel \(N\) of \(q\) is the set of isotropic vectors for \(q\) (Algebra, Chap. IX, §7, no. 1, Cor. of Prop. 2). Let \(a_1, \ldots, a_n\) be the canonical basis of \(\mathbf{R}^n\) . If \(\sum_{i} c_i a_i \in N\) , Lemma 3 shows that we also have \(\sum_{i} |c_i| a_i \in N\) , and hence \(\sum_{i} q_{ji}|c_i| = 0\) for all \(j\) . Let \(I\) be the set of \(i\) such that \(c_i \neq 0\) . If \(j \notin I\) , then \(q_{ji}|c_i| \leqslant 0\) for \(i \in I\) and \(q_{ji}|c_i| = 0\) for \(i \notin I\) , so \(q_{ji} = 0\) for \(j \notin I\) and \(i \in I\) . Assumption \(b)\) thus implies that either \(I = \varnothing\) or \(I = \{1, \ldots, n\}\) . Consequently, every non-zero vector in \(N\) has all its coordinates \(\neq 0\) . If \(\dim N \geqslant 2\) , the intersection of \(N\) with the hyperplane with equation \(x_i = 0\) would be of dimension \(\geqslant 1\) , contrary to what we have just shown. The preceding argument also shows that, if \(\dim N = 1\) , then \(N\) contains a vector all of whose coordinates are \(>0\) . This proves (i).

On the other hand, we know that the eigenvalues of Q are real (Algebra, Chap. IX, §7, no. 3, Prop. 5) and positive since q is positive. Let \(\lambda\) be the smallest of them. The matrix \(Q' = Q - \lambda I_n\) is then the matrix of a degenerate positive form \(q'\) and the off-diagonal elements of \(Q'\) are the same as those of Q. Consequently, \(Q'\) satisfies conditions a), b) and c) of the statement of the lemma. Since the kernel \(N'\) of \(q'\) is the eigenspace of Q corresponding to the eigenvalue \(\lambda\) , assertion (ii) follows from (i).

Lemma 5. Let \(e_1, \ldots, e_n\) be vectors generating \(T\) such that:

\(a) (e_i | e_j) \leqslant 0 \text{ for } i \neq j;\)

\begin{enumerate}
\def\labelenumi{\alph{enumi})}
\setcounter{enumi}{1}
\tightlist
\item
  there does not exist a partition of \(\{1, \ldots, n\}\) into two non-empty subsets I and J such that \((e_i|e_j) = 0\) for \(i \in I\) and \(j \in J\) .
\end{enumerate}

Then there are two possibilities:

\begin{enumerate}
\def\labelenumi{\arabic{enumi})}
\item
  \((e_1, \ldots, e_n)\) is a basis of T;
\item
  \(n = \dim T + 1\) ; there exists a family \((c_i)_{1 \leqslant i \leqslant n}\) of real numbers \(> 0\) such that \(\sum_{i} c_i e_i = 0\) , and any family \((c_i')_{1 \leqslant i \leqslant n}\) of real numbers such that \(\sum_{i} c_i' e_i = 0\) is proportional to \((c_i)_{1 \leqslant i \leqslant n}\) .
\end{enumerate}

Put \(q_{ij} = (e_i|e_j)\) . The matrix \(Q = (q_{ij})\) then satisfies the hypotheses of Lemma 4: conditions a) and b) of Lemma 4 are the same as conditions a) and b) above, and c) is satisfied since \(\sum_{i,j} q_{ij} x_i x_j = \|\sum_{i} x_i e_i\|^2\) . The kernel N of the quadratic form \(q\) on \(\mathbf{R}^n\) , with matrix \(\mathbf{Q}\) , is the set of \((c_1, \ldots, c_n) \in \mathbf{R}^n\) such that \(\sum_{i} c_i e_i = 0\) . If \(N = \{0\}\) , the \(e_i\) are linearly dependent and we are in case 1). If \(\dim N > 0\) , Lemma 4 (i) shows that we are in case 2).

Lemma 6. Let \((e_1, \ldots, e_n)\) be a basis of T such that \((e_i | e_j) \leqslant 0\) for \(i \neq j\) .

\begin{enumerate}
\def\labelenumi{(\roman{enumi})}
\item
  If \(x = \sum_{i} c_i e_i \in \mathbf{T}\) is such that \((x|e_i) \geqslant 0\) for all \(i\) , then \(c_i \geqslant 0\) for all \(i\) .
\item
  If \(x\) and \(y\) are two elements of \(\mathrm{T}\) such that \((x|e_i) \geqslant 0\) and \((y|e_i) \geqslant 0\) for all \(i\) , then \((x|y) \geqslant 0\) . If \((x|e_i) > 0\) and \((y|e_i) > 0\) for all \(i\) , then \((x|y) > 0\) .
\end{enumerate}

Under the hypotheses of (i), assume that \(c_{i} < 0\) for some \(i\) . Let \(f\) be the linear form on \(T\) defined by \(f(e_{i}) = 1\) and

\[
f \left(e _ {j}\right) = - c _ {i} / \left(\sum_ {k = 1} ^ {n} \left| c _ {k} \right|\right) \quad \text { for } j \neq i.
\]

The vectors \(-x, e_1, \ldots, e_n\) then satisfy the hypotheses of Lemma 3 (ii) (taking for \(q\) the metric form on T). We conclude that they are linearly independent, which is absurd. Hence (i). Moreover, if \(x = \sum_{i} c_i e_i \in \mathbf{T}\) and if \(y \in \mathbf{T}\) , then \((x|y) = \sum_{i} c_i (e_i | y)\) , so (ii) follows immediately from (i).

